\section{Subcritical weights and exact Hausdorff dimensions}\label{sec:weights}

\begin{remark}\label{rem:unweighted}
The inequality $\|Q_jx\|\le j\|Q_jx\|$ gives the inclusion from the
weighted zero-set $\{x:j\|Q_jx\|\to0\}=L$ into the unweighted zero-set
$\{x:\|Q_jx\|\to0\}$. To see strictness explicitly, the point $x_0$ constructed
in \texttt{QuarterPointFinal.lean} (with constant digits $d_j\equiv1$) satisfies
$j\|Q_jx_0\|\to1/4$ (\texttt{QuarterPointFinal.\allowbreak lambda\_x0\_eq\_quarter}).
Its unweighted error $\|Q_jx_0\|$ tends to zero, while $x_0\notin L$ by the
weighted classification (Theorem~\ref{thm:class}). The combined witness is
kernel-checked at Layer~A in \texttt{EpiFormal/\allowbreak UnweightedCorrection.lean}
(\texttt{UnweightedCorrection.\allowbreak witness}) using only the three Lean
base axioms.
\end{remark}

Recall that for $0\le\alpha\le1$ we write
$\Lam_\alpha(x)=\limsup_{j\to\infty}j^\alpha\|Q_jx\|$, so that $\Lam_1=\Lam$.

\begin{theorem}[Weighted upper bound, Layer~A in Lean]\label{thm:upper}
For $0\le\alpha\le1$ and every finite $c\in\R$,
$\dimH\{x\in\R:\Lam_\alpha(x)\le c\}\le1-\alpha$.
(Lean: \texttt{WeightedFinal.\allowbreak dimH\_levelSet\ensuremath{\alpha}\_le};
\texttt{\#print axioms} returns \texttt{[propext, Classical.choice, Quot.sound]} only.)
\end{theorem}

\begin{proof}
Fix $0\le\alpha\le1$ and $c\in\R$, and choose $t>\max(c,0)$ (in Lean,
$t=\max(c,0)+1>0$). Put $E_{\alpha,c}:=\{x\in\R:\Lam_\alpha(x)\le c\}$. Since
$Q_j\in\Z$, we have $\|Q_j(x+m)\|=\|Q_jx\|$ for every $m\in\Z$ and $j\ge1$,
hence $\Lam_\alpha(x+m)=\Lam_\alpha(x)$. Thus
\[
E_{\alpha,c}=\bigcup_{m\in\Z}\bigl(m+(E_{\alpha,c}\cap[0,1])\bigr).
\]
By isometry invariance and countable stability of Hausdorff dimension, it
suffices to show that $\dimH(E_{\alpha,c}\cap[0,1])\le1-\alpha$.

For $j\ge1$, set $r_{j,\alpha}:=t/(j^\alpha Q_j)$ and
$A_{j,\alpha}(t):=\{x\in\R:\|Q_jx\|\le t/j^\alpha\}$. Any $x\in E_{\alpha,c}\cap[0,1]$
satisfies $\limsup_{j\to\infty}j^\alpha\|Q_jx\|\le c<t$, so $j^\alpha\|Q_jx\|\le t$
eventually; hence
\[
E_{\alpha,c}\cap[0,1]\subseteq\bigcup_{J=1}^\infty C_{J,\alpha},\qquad
C_{J,\alpha}:=[0,1]\cap\bigcap_{j=J}^\infty A_{j,\alpha}(t).
\]
Fix $J\ge1$. For each $j\ge J$ and $k\in\Z$, let
$I_{j,\alpha}(k):=[k/Q_j-r_{j,\alpha},k/Q_j+r_{j,\alpha}]$ be the closed
interval of center $k/Q_j$ and radius $r_{j,\alpha}$. For any $x\in C_{J,\alpha}$
and $j\ge J$, the nearest integer $k_j(x):=\mathrm{round}(Q_jx)$ belongs to
$\{0,1,\dots,Q_j\}$ and satisfies $|x-k_j(x)/Q_j|=\|Q_jx\|/Q_j\le r_{j,\alpha}$,
so $x\in I_{j,\alpha}(k_j(x))$. At consecutive levels $j$ and $j+1$, the
triangle inequality gives
\begin{equation}\label{eq:weighted-parent-child}
\left|\frac{k_j(x)}{Q_j}-\frac{k_{j+1}(x)}{Q_{j+1}}\right|
\le\left|\frac{k_j(x)}{Q_j}-x\right|+\left|x-\frac{k_{j+1}(x)}{Q_{j+1}}\right|
\le r_{j,\alpha}+r_{j+1,\alpha}=\frac{t}{j^\alpha Q_j}+\frac{t}{(j+1)^\alpha Q_{j+1}}.
\end{equation}
Since $0\le\alpha\le1$ and $j\ge1$, we have $1-\alpha\ge0$ and
$j=j^\alpha j^{1-\alpha}\le j^\alpha(j+1)^{1-\alpha}$, so
$1/j^\alpha\le(j+1)^{1-\alpha}/j$, while
$1/(j+1)^\alpha=(j+1)^{1-\alpha}/(j+1)$. Putting $t'_j:=t(j+1)^{1-\alpha}\ge0$,
we obtain the pointwise radius comparison (\texttt{WeightedFiber.wrad\_le\_rad})
\[
r_{j,\alpha}+r_{j+1,\alpha}
\le\frac{t'_j}{jQ_j}+\frac{t'_j}{(j+1)Q_{j+1}}.
\]
Applying Lemma~\ref{lem:branch} with threshold $t'_j$ shows that for each fixed
parent center $k\in\{0,1,\dots,Q_j\}$, the number of child centers
$k'\in\{0,1,\dots,Q_{j+1}\}$ satisfying \eqref{eq:weighted-parent-child} is at
most
\[
\lfloor16t'_j\rfloor+1\le 16t(j+1)^{1-\alpha}+1\le K(j+1)^{1-\alpha},
\qquad K:=16t+1,
\]
since $(j+1)^{1-\alpha}\ge1$.

Define the surviving center sets $S_{n,\alpha}\subseteq\{0,1,\dots,Q_{J+n}\}$
($n\ge0$) inductively by $S_{0,\alpha}:=\{0,1,\dots,Q_J\}$ and letting
$S_{n+1,\alpha}$ consist of all $k'\in\{0,1,\dots,Q_{J+n+1}\}$ within
$r_{J+n,\alpha}+r_{J+n+1,\alpha}$ of some $k/Q_{J+n}$ with $k\in S_{n,\alpha}$.
By induction on $n\ge0$, $C_{J,\alpha}\subseteq\bigcup_{k\in S_{n,\alpha}}I_{J+n,\alpha}(k)$
and the cardinality $|S_{n,\alpha}|$ satisfies
(\texttt{WeightedDimH.card\_wsurvSet\_le})
\[
|S_{n,\alpha}|
\le(Q_J+1)\prod_{i=0}^{n-1}\bigl(K(J+i+1)^{1-\alpha}\bigr)
\le(Q_J+1)K^n\bigl((J+n)!\bigr)^{1-\alpha},
\]
because $\prod_{i=0}^{n-1}(J+i+1)=(J+n)!/J!\le(J+n)!$. By
Lemma~\ref{lem:factorial}, $Q_{J+n}\ge(J+n)!$, and $(J+n)^\alpha\ge1$, so each
interval $I_{J+n,\alpha}(k)$ has diameter
\[
\operatorname{diam}I_{J+n,\alpha}(k)=2r_{J+n,\alpha}
=\frac{2t}{(J+n)^\alpha Q_{J+n}}\le\frac{2t}{(J+n)!}\longrightarrow0
\qquad(n\to\infty).
\]
Now fix any exponent $s>1-\alpha$, and put $\beta:=s-(1-\alpha)>0$. For every
$\delta>0$ and all sufficiently large $n$, $\{I_{J+n,\alpha}(k):k\in S_{n,\alpha}\}$
is a finite $\delta$-cover of $C_{J,\alpha}$ whose $s$-dimensional Hausdorff
covering cost satisfies
\[
\sum_{k\in S_{n,\alpha}}\bigl(\operatorname{diam}I_{J+n,\alpha}(k)\bigr)^s
\le |S_{n,\alpha}|\left(\frac{2t}{(J+n)!}\right)^s
\le(Q_J+1)(2t)^s K^{-J}\cdot K^{J+n}\bigl((J+n)!\bigr)^{1-\alpha-s}.
\]
With $m=J+n$ and $K^{1/\beta}>0$, the factor $K^m(m!)^{1-\alpha-s}=K^m(m!)^{-\beta}=\bigl((K^{1/\beta})^m/m!\bigr)^\beta$
satisfies, via the elementary factorial lower bound $m!\ge(m/\ee)^m$,
\[
0\le K^m(m!)^{-\beta}=\left(\frac{(K^{1/\beta})^m}{m!}\right)^\beta
\le\left(\frac{\ee K^{1/\beta}}{m}\right)^{\beta m}\longrightarrow0
\qquad(m=J+n\to\infty)
\]
(formalized in \texttt{WeightedPhase.\allowbreak content\_\allowbreak core\_\allowbreak tendsto} via
\texttt{FloorSemiring.\allowbreak tendsto\_\allowbreak pow\_\allowbreak div\_\allowbreak factorial\_\allowbreak atTop} and continuity of
$u\mapsto u^\beta$ at $0$). Hence $\mathcal H_\delta^s(C_{J,\alpha})=0$ for all
$\delta>0$, so $\mathcal H^s(C_{J,\alpha})=0$ and $\dimH C_{J,\alpha}\le s$.
Since $s>1-\alpha$ was arbitrary, letting $s\downarrow1-\alpha$ gives
$\dimH C_{J,\alpha}\le1-\alpha$ for every $J\ge1$, and countable stability over
$J\ge1$ and $m\in\Z$ yields $\dimH E_{\alpha,c}\le1-\alpha$.
\end{proof}

\paragraph{Lean status.} Theorem~\ref{thm:upper} is kernel-checked at Layer~A
as \texttt{WeightedFinal.\allowbreak dimH\_levelSet\ensuremath{\alpha}\_le},
assembled through \texttt{WeightedFiber.lean}, \texttt{WeightedDimHCore.lean},
\texttt{WeightedUpper.lean}, and \texttt{WeightedFinal.lean}. Its axiom audit
reports only the three base axioms.

\begin{theorem}[Prescribed limits at subcritical weights, Layer~A in Lean]\label{conj:weight}
Let $0\le\alpha<1$. For $\alpha>0$ take any $c\ge0$; for $\alpha=0$
take $0\le c\le1/2$. Then
\[
 \dimH\{x\in[0,1]: j^\alpha\|Q_jx\|\longrightarrow c\}=1-\alpha .
\]
In particular $\dimH\{x\in\R:\Lam_\alpha(x)\le c\}=1-\alpha$ for every
finite $c\ge0$, including $c=0$.
(Lean, module \texttt{WeightedExactDim}:
\texttt{dimH\_\allowbreak exactLimitSet\ensuremath{\alpha}\_eq},
\texttt{dimH\_\allowbreak levelSet\ensuremath{\alpha}\_eq}, and
\texttt{dimH\_\allowbreak apprExponent\ensuremath{\alpha}\_eq\_\allowbreak zero};
\texttt{\#print axioms} returns \texttt{[propext, Classical.choice, Quot.sound]} only.)
\end{theorem}

\begin{proof}
Write $E_{\alpha}^{\mathrm{lim}}(c):=\{x\in[0,1]:j^\alpha\|Q_jx\|\to c\}$. Since
$j^\alpha\|Q_jx\|\to c$ implies $\Lam_\alpha(x)=\limsup_{j\to\infty}j^\alpha\|Q_jx\|=c$,
we have $E_{\alpha}^{\mathrm{lim}}(c)\subseteq\{x\in\R:\Lam_\alpha(x)\le c\}$, so
Theorem~\ref{thm:upper} immediately gives the upper bound
$\dimH E_{\alpha}^{\mathrm{lim}}(c)\le1-\alpha$.

We now give the complete paper proof of the lower bound
$\dimH E_{\alpha}^{\mathrm{lim}}(c)\ge1-\alpha$. For $j\ge1$, put
\[
w_j:=\frac{j^{-\alpha}}{\log(j+2)},\qquad r_j:=\frac{w_j}{Q_j}.
\]
Choose an integer $N\ge3$ large enough that $w_j<1/8$ for all $j\ge N$, and
also $cj^{-\alpha}+w_j\le1/2$ for all $j\ge N$ whenever either ($\alpha>0$,
$c\ge0$) or ($\alpha=0$, $0\le c<1/2$). Define the target fractional shift
$\theta_j\in[w_j,1/2-w_j]$ for $j\ge N$ by clamping $cj^{-\alpha}$ away from the
half-integer boundary:
\[
\theta_j:=\begin{cases}
w_j,&\text{if }c=0,\\[2pt]
cj^{-\alpha},&\text{if }c>0\text{ and }(\alpha>0\text{ or }\alpha=0,c<1/2),\\[2pt]
\tfrac12-w_j,&\text{if }\alpha=0\text{ and }c=\tfrac12.
\end{cases}
\]
(Note why the boundary adjustment $\theta_j=1/2-w_j$ is needed in the endpoint
case $\alpha=0,c=1/2$, and $\theta_j=w_j$ when $c=0$: if one centered directly
at $(k+1/2)/Q_j$, points on the right half of the interval would cross the
half-integer $k+1/2$ where $y\mapsto\|y\|$ folds back, whereas with center
$(k+\theta_j)/Q_j$ and radius $r_j=w_j/Q_j$, every point $x$ in the closed
interval $I_j(k):=[(k+\theta_j-w_j)/Q_j,(k+\theta_j+w_j)/Q_j]$ satisfies
$k\le k+\theta_j-w_j\le Q_jx\le k+\theta_j+w_j\le k+1/2$. On $[k,k+1/2]$ the
distance to the nearest integer is given without folding by
$\|Q_jx\|=Q_jx-k\in[\theta_j-w_j,\theta_j+w_j]$.) In particular, in the case
$\alpha=0,c=1/2$, for $x\in I_j(k)$ we have
$Q_jx\in[k+1/2-2w_j,k+1/2]$, so
\[
\|Q_jx\|=\frac12-\left(k+\frac12-Q_jx\right)\in\left[\frac12-2w_j,\frac12\right]
\quad\Longrightarrow\quad
\bigl|\|Q_jx\|-\tfrac12\bigr|\le2w_j=\frac{2}{\log(j+2)}\longrightarrow0,
\]
and in all three cases every $x\in I_j(k)$ satisfies
\[
\bigl|j^\alpha\|Q_jx\|-c\bigr|
\le j^\alpha\bigl|\|Q_jx\|-\theta_j\bigr|+\bigl|j^\alpha\theta_j-c\bigr|
\le j^\alpha w_j+j^\alpha w_j=\frac{2}{\log(j+2)}\longrightarrow0
\qquad(j\to\infty).
\]
Moreover, for fixed $j\ge N$ and distinct integers $k_1\ne k_2$, the distance
between the centers of $I_j(k_1)$ and $I_j(k_2)$ is
$|k_1-k_2|/Q_j\ge1/Q_j>2w_j/Q_j=2r_j$ (since $w_j<1/8<1/2$), so the level-$j$
intervals $\{I_j(k):k\in\Z\}$ are pairwise disjoint.

Next we bound the number of strictly nested level-$(j+1)$ intervals inside a
given parent interval $I_j(k)$ ($j\ge N$). A child interval
$I_{j+1}(k')=[c_{j+1}(k')-r_{j+1},c_{j+1}(k')+r_{j+1}]$ with center
$c_{j+1}(k')=(k'+\theta_{j+1})/Q_{j+1}$ is contained in
$I_j(k)=[c_j(k)-r_j,c_j(k)+r_j]$ ($c_j(k)=(k+\theta_j)/Q_j$) if and only if
$|c_{j+1}(k')-c_j(k)|\le r_j-r_{j+1}$, which is equivalent to requiring the
integer $k'$ to lie in the closed interval $[A_j(k),B_j(k)]$ where
\[
A_j(k):=Q_{j+1}\bigl(c_j(k)-(r_j-r_{j+1})\bigr)-\theta_{j+1},\qquad
B_j(k):=Q_{j+1}\bigl(c_j(k)+(r_j-r_{j+1})\bigr)-\theta_{j+1}.
\]
The length of $[A_j(k),B_j(k)]$ is independent of the parent index $k$ and of
the shifts $\theta_j,\theta_{j+1}$:
\[
W_j:=B_j(k)-A_j(k)=2(r_j-r_{j+1})Q_{j+1}
=2\left(w_j\frac{Q_{j+1}}{Q_j}-w_{j+1}\right).
\]
Since any closed real interval $[A,B]$ of length $W=B-A\ge1$ contains at least
$\lfloor B\rfloor-\lceil A\rceil+1>W-1$ integers, the number of eligible child
centers $k'\in\Z$ with $I_{j+1}(k')\subseteq I_j(k)$ is at least
$2(r_j-r_{j+1})Q_{j+1}-1$. Using $Q_{j+1}=(4j+2)Q_j+Q_{j-1}>(4j+2)Q_j$ from
\eqref{eq:rec} and $w_{j+1}\le w_j$, we have
\[
2(r_j-r_{j+1})Q_{j+1}-1
=2w_j\frac{Q_{j+1}}{Q_j}-2w_{j+1}-1
\ge 2w_j(4j+2)-2w_j-1
=\frac{8j^{1-\alpha}+2j^{-\alpha}}{\log(j+2)}-1.
\]
Since $0\le\alpha<1$, $j^{1-\alpha}/\log(j+2)\to\infty$ as $j\to\infty$.
Increasing $N$ if necessary so that $7j^{1-\alpha}/\log(j+2)\ge1$ and
$b_j:=\lfloor j^{1-\alpha}/\log(j+2)\rfloor\ge2$ for all $j\ge N$, we obtain
\[
2(r_j-r_{j+1})Q_{j+1}-1\ge\frac{j^{1-\alpha}}{\log(j+2)}\ge b_j\ge2
\qquad\text{for all }j\ge N,
\]
so every parent interval $I_j(k)$ ($j\ge N$) contains at least $b_j$ pairwise
disjoint child intervals $I_{j+1}(k')$.

Choose $N$ large enough that $I_N(1)\subseteq(0,1)$. Construct a homogeneous
Cantor tree by starting at level $N$ with the single root interval $I_N(1)$
($\mathcal E_N:=\{1\}$), and for each $j\ge N$ and each retained parent
$k\in\mathcal E_j$, keeping exactly $b_j$ child indices $k'\in\mathcal E_{j+1}$
with $I_{j+1}(k')\subseteq I_j(k)$ (for instance, the first $b_j$ integers
$\lceil A_j(k)\rceil+m$, $0\le m<b_j$). Then at each level $n\ge N$ there are
exactly
\[
M_n:=\prod_{j=N}^{n-1}b_j
\]
pairwise disjoint closed intervals $\{I_n(k):k\in\mathcal E_n\}$ of radius
$r_n=w_n/Q_n$, and their compact intersection
\[
K:=\bigcap_{n=N}^\infty\bigcup_{k\in\mathcal E_n}I_n(k)\subseteq[0,1]
\]
is nonempty and satisfies $K\subseteq E_{\alpha}^{\mathrm{lim}}(c)$. Let $\mu$
be the unique Borel probability measure on $K$ assigning equal mass
$\mu(I_n(k))=M_n^{-1}$ to each of the $M_n$ retained level-$n$ intervals ($n\ge N$).

To estimate the asymptotic growth of $\log M_n$ and $\log Q_{n+1}$, observe
that for $j\ge N$,
\[
\frac{j^{1-\alpha}}{2\log(j+2)}\le\frac{j^{1-\alpha}}{\log(j+2)}-1
<b_j\le\frac{j^{1-\alpha}}{\log(j+2)},
\]
so $\log b_j=(1-\alpha)\log j-\log\log(j+2)+O(1)$. By integral comparison for
the monotone function $t\mapsto\log t$,
$\sum_{j=N}^{n-1}\log j=\int_1^n\log t\,dt+O(\log n)=n\log n-n+O(\log n)$, and
$\sum_{j=N}^{n-1}\log\log(j+2)\le n\log\log(n+2)$, which gives
\begin{equation}\label{eq:log-Mn-asymp}
\log M_n=\sum_{j=N}^{n-1}\log b_j
=(1-\alpha)n\log n+O(n+n\log\log n)=(1-\alpha)n\log n+o(n\log n).
\end{equation}
Similarly, from $4jQ_j<(4j+2)Q_j<Q_{j+1}\le(4j+3)Q_j$ for all $j\ge1$, taking
logarithms and summing from $j=1$ to $n$ gives
\[
\log Q_{n+1}=\sum_{j=1}^n\log\frac{Q_{j+1}}{Q_j}
=\sum_{j=1}^n\bigl(\log j+\log 4+O(1/j)\bigr)=n\log n+O(n).
\]

Now fix any $0<s<1-\alpha$. By \eqref{eq:log-Mn-asymp} and
$\log Q_{n+1}=n\log n+O(n)$, we have
\begin{align*}
\log\left(\frac{3}{M_n Q_{n+1}^{-s}}\right)
&=s\log Q_{n+1}-\log M_n+\log 3\\
&=-\bigl((1-\alpha)-s\bigr)n\log n+o(n\log n)\longrightarrow-\infty
\qquad(n\to\infty).
\end{align*}
Hence we may choose an index $N_s\ge N$ such that
\begin{equation}\label{eq:mass-scale-bound}
\frac{3}{M_n}\le Q_{n+1}^{-s}\qquad\text{for all }n\ge N_s.
\end{equation}
Let $\{U_i\}_{i\ge1}$ be any countable cover of $K$ by nonempty sets of diameter
$0<h_i:=\operatorname{diam}(U_i)<Q_{N_s}^{-1}$ (if $\operatorname{diam}(U_i)=0$,
then $U_i$ is a singleton and $\mu(U_i)=0$ since $M_n^{-1}\to0$). For each $i$,
since $Q_n$ is strictly increasing to $\infty$, there exists a unique integer
$n_i\ge N_s$ such that
\[
Q_{n_i+1}^{-1}\le h_i<Q_{n_i}^{-1}.
\]
The set $U_i$ is contained in a closed interval $[x_0,x_0+h_i]$ of length
$h_i<Q_{n_i}^{-1}$. If a level-$n_i$ interval $I_{n_i}(k)$ ($k\in\Z$) intersects
$[x_0,x_0+h_i]$, then its center $c_{n_i}(k)=(k+\theta_{n_i})/Q_{n_i}$ must lie
in $[x_0-r_{n_i},x_0+h_i+r_{n_i}]$, which has length
$h_i+2r_{n_i}<Q_{n_i}^{-1}+2w_{n_i}Q_{n_i}^{-1}<(1+1/4)Q_{n_i}^{-1}<2Q_{n_i}^{-1}$.
Since the centers $c_{n_i}(k)$ ($k\in\Z$) form a grid of spacing exactly
$Q_{n_i}^{-1}$, an interval of length $<2Q_{n_i}^{-1}$ can contain at most $3$
grid points $c_{n_i}(k)$ (if it contained $4$ points $k_1<k_2<k_3<k_4$, its
length would be at least $(k_4-k_1)Q_{n_i}^{-1}\ge3Q_{n_i}^{-1}$). Thus $U_i$
intersects at most $3$ retained level-$n_i$ intervals $I_{n_i}(k)$
($k\in\mathcal E_{n_i}$). Because $\mu$ is supported on
$\bigcup_{k\in\mathcal E_{n_i}}I_{n_i}(k)$ and each $I_{n_i}(k)$ has measure
$M_{n_i}^{-1}$, it follows from \eqref{eq:mass-scale-bound} and
$Q_{n_i+1}^{-1}\le h_i$ that
\[
\mu(U_i)\le\frac{3}{M_{n_i}}\le Q_{n_i+1}^{-s}\le h_i^s=\bigl(\operatorname{diam}(U_i)\bigr)^s.
\]
Summing over the cover $\{U_i\}$ yields
\[
\sum_{i}\bigl(\operatorname{diam}(U_i)\bigr)^s
\ge\sum_{i}\frac{3}{M_{n_i}}\ge\sum_{i}\mu(U_i)\ge\mu(K)=1.
\]
Taking the infimum over all countable $\delta$-covers with
$0<\delta<Q_{N_s}^{-1}$ and letting $\delta\downarrow0$ gives
$\mathcal H^s(E_{\alpha}^{\mathrm{lim}}(c))\ge\mathcal H^s(K)\ge1>0$, and therefore
$\dimH E_{\alpha}^{\mathrm{lim}}(c)\ge s$. Letting $s\uparrow1-\alpha$ proves
$\dimH E_{\alpha}^{\mathrm{lim}}(c)\ge1-\alpha$, establishing the exact identity
$\dimH E_{\alpha}^{\mathrm{lim}}(c)=1-\alpha$.

Finally, for any finite $c\ge0$ (including $c=0$), we have the inclusions
\[
E_{\alpha}^{\mathrm{lim}}(0)
\subseteq\{x\in\R:\Lam_\alpha(x)=0\}
\subseteq\{x\in\R:\Lam_\alpha(x)\le c\}.
\]
The lower bound $\dimH E_{\alpha}^{\mathrm{lim}}(0)=1-\alpha$ just proved and
the upper bound $\dimH\{x\in\R:\Lam_\alpha(x)\le c\}\le1-\alpha$ from
Theorem~\ref{thm:upper} squeeze both sets to have exact Hausdorff dimension
$1-\alpha$.

The Lean formalization in \texttt{EpiFormal/\allowbreak WeightedExactDim.lean}
certifies the same lower bound at Layer~A for every $0<s<1-\alpha$ by taking
polynomial window widths $w_j=j^{-\alpha-\varepsilon}$ with
$\varepsilon=(1-\alpha-s)/2>0$, clamping the target shift into
$[w_j,1/2-w_j]$, and applying compactness plus finite-tree counting directly to
the definition of Hausdorff measure to establish
$\mathcal H^s(K)\ge1$ (\texttt{one\_le\_hausdorffMeasure\_cantorSet}).
\end{proof}

For $\alpha=0$ the dimension assertion also follows from Corollary~8 of
\cite{BLR} with constant target $c$, by the distance inequality used
above. For $0<\alpha<1$, condition~(1.6) of their Theorem~7 fails because
$\log Q_{j+1}/\log Q_j\to1$. The lower bound in the paper proof above can also
be read from the homogeneous Moran bounds of Feng, Wen and Wu \cite{FWW}: the
retained level has $\log M_n\sim(1-\alpha)n\log n$ and both bounding logarithmic
length scales are asymptotic to $n\log n$. The independent covering upper bound
(Theorem~\ref{thm:upper}) applies to the full prescribed-limit set and sublevel
set.

